\section{Discussion}
\label{sec:discussion}

\noindent\textbf{Why executable code is the key enabling representation.}
The core claim of this paper is not about LLMs, CadQuery, or any specific provider.
It is about the choice of representation.
Mesh generation and point-cloud generation do not admit exact geometric verification:
there is no clear threshold at which a generated mesh is ``correct.''
Code generation does: a program either executes and matches the target, or it does not.
This distinction is what makes our entire pipeline possible---verification is free, fast,
and precise.
For any engineering domain where the ground truth can be represented as an executable
artifact, the same approach applies.

\noindent\textbf{Verification as curation, not evaluation.}
Standard practice in ML is to use geometric metrics as post-hoc evaluation on a held-out
test set.
We use IoU~$\geq 0.99$ as a real-time curation gate during training data construction.
This is a stronger and more unusual role: the verifier must run in seconds (not minutes),
must be deterministic (not model-based), and must be precise enough that near-correct
programs are rejected rather than accepted.
Our OCCT-based verifier satisfies all three: ${\approx}2$ seconds per pair, deterministic
Boolean intersection, and rejection sensitivity down to ${\approx}5\,\text{mm}^3$ volume
discrepancy on typical parts.

\noindent\textbf{Why skill-based modularity matters for hard engineering data.}
The 134 manually fixed parts in our dataset would not exist under a rigid automated
pipeline---they represent cases where no provider in the cascade achieved
IoU~$\geq 0.97$, but the correct CadQuery program exists and is findable by a skilled
human using the repair skill.
The key enabling factor is that ``manual checking'' is a first-class component of the
system, not an afterthought.
The skill framework makes it easy to route hard cases to human inspection, collect the
repaired programs, and re-integrate them into the verified store with full provenance.
This is the engineering advantage of modularity: the hardest cases get the most expensive
resource (human time), but only those cases.

\noindent\textbf{Limitations.}
The primary limitation of this work is the absence of downstream retraining results.
Whether the 1,784 verified pairs improve a CAD reconstruction model over existing
datasets---\eg, over training on the raw Fusion360 Gallery with weak supervision---remains
to be shown.
A second limitation is that geometric IoU~$\geq 0.99$ is necessary but not sufficient for
training quality: two STEP files with identical geometry may have different construction
histories (\eg, a feature built by join vs.\ by cut), and a downstream model trained on
the wrong construction history may generalize poorly.
Addressing this would require construction-history verification beyond solid geometry,
which is an open problem.
Third, the 109 near-misses currently in the fix queue represent failure modes that have
resisted both automated repair and our developed heuristics; a more systematic analysis of
these remaining cases is future work.

\section{Conclusion}
\label{sec:conclusion}

We presented a skill-based, verification-guided framework for CAD data generation.
By reformulating complex CAD generation as verifiable program synthesis---executable
programs that can be geometrically validated against ground truth---and by organizing the
system around reusable skills rather than a rigid workflow, we convert noisy LLM outputs
into high-precision, training-ready supervision without human annotation.
The result is a verified corpus of 1,784 program--geometry pairs covering complex
multi-operation industrial parts, all with IoU~$\geq 0.99$ against ground truth.
Our analysis shows that 92.5\% of verified pairs are produced automatically by the
provider cascade, while 7.5\% require targeted manual repair---a small but important tail
that captures genuinely hard cases the automated pipeline cannot reach.
More broadly, our results suggest that for complex engineering tasks, the combination of
executable representations, strict geometric verification, and modular skill composition
provides a practical path to scalable data curation that does not require expensive human
labeling.
